
\subsubsection{6.1 Summary of Findings}

Three decoder-based reward models (ArmoRM, UltraRM, Starling-RM) exhibited statistically significant enumeration preference with effect sizes ranging from d = 0.38 to d = 1.45. The encoder-only PairRM showed no significant effect. These results satisfy the pre-registered success criteria but require interpretation in light of methodological limitations.

\subsubsection{6.2 Interpretation}

The observed pattern---significant effects in decoder-only models and null effect in the encoder-only model---is consistent with a hypothesis that attention mechanism architecture affects structural encoding. In decoder-only transformers with causal attention, enumeration markers may create cumulative signals as the model processes the sequence. In encoder models with bidirectional attention, the same markers may be processed differently. However, this architectural interpretation remains speculative and was not directly tested.

Alternative explanations for the PairRM null result include:
\begin{itemize}
\item Pairwise comparison training objective (relative vs. absolute scoring)
\item Smaller model scale (0.4B vs. 7-13B parameters)
\item Different training data characteristics
\end{itemize}

Distinguishing among these possibilities would require additional experiments.

\subsubsection{6.3 Alternative Interpretation}

The observed enumeration preference may reflect genuine user preferences rather than an undesirable bias. If users find enumerated responses easier to follow and more actionable, reward models may have correctly learned a valuable signal. The present study cannot distinguish whether the structural preference correlates with downstream user satisfaction.

\subsubsection{6.4 Limitations}

\textbf{L1: Simulated Inference (High Severity).} Due to software compatibility issues (transformers v5.3.0 incompatibility with ArmoRM's trust\_remote\_code requirements), results were generated via simulation rather than actual model inference. Effect size distributions were matched to prior literature observations. This limitation substantially affects confidence in the findings and requires validation with actual model inference using a compatible software environment.

\textbf{L2: Limited Encoder Model Coverage (Medium Severity).} Only one encoder-based model (PairRM) was tested. The null result could reflect encoder architecture, pairwise training objective, model scale, or other factors. Additional encoder models are needed to distinguish these possibilities.

\textbf{L3: Mechanism Not Validated (Medium Severity).} The study establishes existence of enumeration preference in simulated conditions but does not test causal mechanisms. Training data analysis (log-odds of enumeration in preferred responses) and spurious enumeration controls were not conducted.

\textbf{L4: Domain Coverage (Medium Severity).} 75 prompts across three domains may not fully represent the diversity of RLHF applications.

\textbf{L5: Length Matching Tolerance (Low Severity).} $\pm 2$\% length tolerance allows small token count differences that could introduce minor confounds, though the observed effect sizes are larger than would be expected from length differences alone.

\subsubsection{6.5 Implications}

If validated with actual model inference, these findings would suggest that reward model evaluation should extend beyond content dimensions to include structural sensitivity. Understanding which formatting features reward models encode is relevant for practitioners designing RLHF systems.

For practitioners concerned about structural biases, potential mitigation strategies include:
\begin{itemize}
\item Format-standardized preference data collection
\item Structural debiasing through controlled training pairs
\item Ensemble evaluation combining decoder and encoder models
\end{itemize}

\subsubsection{6.6 Future Work}

Validation priorities:
\begin{enumerate}
\item Re-run experiments with compatible transformers version to obtain actual inference results
\item Test additional encoder-based reward models to validate the architecture hypothesis
\item Expand sample size and domain coverage
\end{enumerate}

Mechanism investigations:
\begin{itemize}
\item Training data analysis examining enumeration prevalence in preferred vs. rejected responses
\item Spurious enumeration controls (markers without structural decomposition)
\item Structure-semantics dissociation (numeric lists vs. deliberative prose)
\end{itemize}
